\documentclass[11pt]{article}

\usepackage[margin=1.05in]{geometry}
\usepackage{amsmath,amssymb,amsthm}
\usepackage{booktabs}
\usepackage{enumitem}
\usepackage{xcolor}
\usepackage[colorlinks=true,linkcolor=blue!45!black,citecolor=blue!45!black,urlcolor=blue!45!black]{hyperref}

\newcommand{\R}{\mathbb{R}}
\newcommand{\E}{\mathbb{E}}
\newcommand{\tr}{\operatorname{tr}}
\newcommand{\code}[1]{\texttt{\small #1}}
\newcommand{\Var}{\operatorname{Var}}
\newcommand{\khat}{\widehat{k}}
\newcommand{\Rhat}{\widehat{R}}

\theoremstyle{remark}
\newtheorem{remark}{Remark}

\title{\bfseries Three obstructions to Stein variational gradient descent\\[2pt]
on a MAGI posterior}
\author{}
\date{\today}

\begin{document}
\maketitle

\begin{abstract}
\noindent
Stein variational gradient descent (SVGD) is an attractive candidate for MAGI \cite{yang2021}
posteriors: it is deterministic, it parallelises trivially, and it needs only the score. We spent
four investigations on it and recommend against it. This report sets out why, because the reasons
are more interesting than the verdict and two of them are not specific to MAGI.

There are three obstructions, and each would be sufficient alone. First, the variance collapse is
dimensional. On these problems $d$ runs from $106$ to $608$, and the equilibrium variance ratio
under the median-heuristic bandwidth is $\ln K/d$ --- measured to three decimals over $d \in
[50, 608]$ and $K \in [10, 3200]$ --- so correct dispersion needs $K \sim e^{d}$. Ba et
al.\ \cite{ba2022} prove the $\Theta(1/d)$ scaling; the $\ln K$ numerator is an artifact of the
$1/\log K$ factor that implementations following Liu and Wang \cite{liu2016} apply, and dropping it
is worth a factor of $0.582K/\ln K$, or $39$ at $K = 400$. Second, that obstruction \emph{is}
removable: at a large fixed bandwidth in coordinates whitened by the Laplace metric, SVGD has a
genuine attractor, the residual error is a pure scale deficit, and the Stein identity measures that
deficit without a reference --- rescaling by $1/\sqrt{R}$ brings the ensemble to $0.57$ to $0.96$
times the reference chain's own agreement floor. Third, and fatally, the fixed point so obtained is
displaced in the \emph{parameters}, along the axis from the joint mode to the posterior mean, by
that axis's own length. It is exact for a Gaussian target and biased for anything else. That
displacement is precisely the correction the deterministic pipeline exists to compute, so SVGD
returns the quantity being estimated.

We also report a bug of our own that invalidated an earlier verdict on preconditioning, two of our
own hypotheses that measurement overturned, and the three things from this work worth keeping: the
$1/\log K$ deletion, the Stein dispersion ratio as a diagnostic, and the design rule
$K \gtrsim e^{p}$ for SVGD on a $p$-dimensional marginal.
\end{abstract}

\section{What would count as success}
\label{sec:setting}

MAGI places a Gaussian process prior on each state component of $\dot x = f(x,\theta)$ and
constrains it to satisfy the dynamics on a discretisation grid. The unknown is
$x = (\theta, X) \in \R^{d}$ with $d = p + nD$; across the four systems we use, $d$ runs from
$106$ to $608$ and $p$ from $3$ to $7$. Almost all of the dimension is state.

Three of the four systems have a usable reference: $32$ preconditioned NUTS chains, $96{,}000$
draws, $\Rhat \le 1.007$. Hes1 does not --- $\Rhat = 1.76$ with $13\%$ divergences --- and is
excluded from every quantitative statement here, because without a reference nothing measured on it
can be scored.

Scoring is by energy distance in Mahalanobis coordinates, always reported against its own
\textbf{floor}: the same statistic between two independent subsamples of the reference at the
ensemble's particle count. A value of $1.0\times$ the floor means the ensemble is as close to the
reference as the reference is to itself, which is the most a finite reference can establish.
Marginal standard deviations are \emph{not} used as a score, and this is the specific trap the
problem sets: an early run held every marginal sd at $0.991$ of the reference --- visually perfect
credible intervals --- while sitting $45\times$ the energy floor from the target. SVGD's collapse
is anisotropic and no marginal sees it.

The incumbent is the deterministic pipeline of \cite{companion2}: profile the states out by Laplace
at each $\theta$, integrate the $p$-dimensional remainder by importance sampling. It reaches the
reference's own agreement floor on all three scorable systems --- largest parameter error $0.0126$
against a floor of $0.0100$, $0.0093$ against $0.0081$, $0.0119$ against $0.0405$ --- in four to
six seconds. Any case for SVGD has to be made against that.

\section{Obstruction 1: the collapse is dimensional}
\label{sec:dim}

\subsection{The original observation, and the control that explains it}

Started \emph{at} a correct posterior sample on FitzHugh--Nagumo and run, SVGD moves away from it:
the energy distance degrades from $0.085$ to $6.89$ under the standard RBF kernel and $3.80$ under
the density-reweighted kernel of \cite{huang2023}, and the Stein-identity ratio falls from $1$ to
$0.019$ and $0.24$. Since the starting point was correct, the motion is unambiguously away from it.

We re-ran that after correcting a Gaussian-process hyperparameter fit and changing the benchmark
integrator from forward Euler to RK4, expecting the numbers to move. They do not: the standard
kernel reproduces to three significant figures (energy $6.895 \to 6.847$, Stein $0.0190 \to
0.0194$), and the explicit $2\times2$ of $\{$RK4, Euler$\}\times\{$fixed, old$\}$ hyperparameters
gives Stein $0.0192$ to $0.0194$ in every cell. Neither correction has any bearing on it.

The reason they have no bearing is that the effect is not a property of the MAGI posterior at all.
The control that settles it is embarrassingly simple and had not been run: SVGD on
$\mathcal N(0, I_{325})$, started at exact draws. It collapses to a flat $1.8\%$ of the correct
variance --- $89\times$ the floor, the same as on the real posterior. No MAGI, no anisotropy, no
non-Gaussianity.

\subsection{The law}

Sweeping $d$ and $K$ on isotropic Gaussian targets, started at exact draws so there is no burn-in
to mistake for an equilibrium, the mean variance ratio times $d$ collapses onto a function of $K$
alone:

\begin{table}[t]
\centering\small
\begin{tabular}{lrrrr}
\toprule
$d \backslash K$ & 50 & 100 & 400 & 1600\\
\midrule
50   & 3.915 & 4.512 & 5.452 & 8.602\\
100  & 3.912 & 4.605 & 5.774 & 7.924\\
200  & 3.912 & 4.582 & 5.914 & 7.756\\
325  & 3.917 & 4.601 & 5.956 & 8.072\\
608  & 3.890 & 4.605 & 5.978 & 7.469\\
\midrule
$\ln K$ & \textbf{3.912} & \textbf{4.605} & \textbf{5.991} & \textbf{7.378}\\
\bottomrule
\end{tabular}
\caption{$d \cdot \Var_{\mathrm{SVGD}}/\Var_{\mathrm{target}}$ against $\ln K$, standard RBF kernel
under the median heuristic with the $1/\log K$ factor. Prodigy, $2000$ iterations, started at exact
draws.}
\label{tab:law}
\end{table}

\begin{equation}
\label{eq:law}
  \frac{\Var_{\mathrm{SVGD}}}{\Var_{\mathrm{target}}} \;=\; \frac{\ln K}{d},
\end{equation}
to three decimals over an order of magnitude in $d$, confirmed separately by sweeping $K$ from $10$
to $3200$ at fixed $d$. At $d = 325$ and $K = 400$ that is $1.8\%$.

The mechanism is a feedback loop in the bandwidth rule. The median heuristic measures $h$ from the
\emph{ensemble}, so contraction tightens the bandwidth, which weakens the repulsion at the scale
that matters, which permits further contraction. The equilibrium is where that loop closes, and it
closes at a normalised bandwidth of $h \approx 2$ regardless of where it started.

\subsection{What is ours and what is not}

The $1/d$ scaling is not ours. Ba, Erdogdu, Ghassemi, Sun, Suzuki, Wu and Zhang \cite{ba2022}
prove, for an isotropic Gaussian target in the proportional limit $d/K \to \gamma > 1$ under the
\emph{plain} median heuristic $h = \mathrm{Med}$, that
$\Var_{\mathrm{SVGD}}/\Var = (e-1)^{-1}K/d = 0.582\,K/d$ (their Corollary 4), conditional on a
near-orthogonality assumption they verify numerically rather than prove. Section~\ref{sec:dim}'s
\eqref{eq:law} is the same $1/d$ with a different numerator, and the difference is entirely the
bandwidth rule:

\begin{table}[t]
\centering\small
\begin{tabular}{llll}
\toprule
& bandwidth & law & at $K = 400$, $d = 800$\\
\midrule
Ba et al.\ \cite{ba2022}, Cor.~4 & $h = \mathrm{Med}$ & $0.582\,K/d$ & $0.291$\\
this work, and what implementations ship & $h = \mathrm{Med}/\ln K$ & $\ln K/d$ & $0.0075$\\
\bottomrule
\end{tabular}
\end{table}

Running both conventions through one harness reproduces both laws to within $1\%$ --- measured
$0.29026$ against a predicted $0.29099$, and $0.00750$ against $0.00749$ --- which validates the
harness against published theory and isolates a fact worth acting on. Liu and Wang \cite{liu2016}
recommend $h = \mathrm{Med}/\log K$, and that is what our library used. The cost is
$0.582K/\ln K$: a factor of $39$ at $K = 400$ and $281$ at $K = 4000$, always against the log.

Deleting it is one line. On the three real posteriors at $K = 400$ it improves the energy distance
by $1.8$ to $2.4\times$ and the Stein ratio by $5$ to $9\times$, better in all six cells tested; on
$\mathcal N(0, I_{20})$ at $K = 200$ the variance ratio goes from $0.265$ to $0.929$ against a
target of $1$. It is free and it is a strict improvement, and we have made the change. It is not a
fix: both variants are $O(1/d)$, and \eqref{eq:law} rearranged says correct dispersion needs
$K \sim e^{d}$. At $d = 325$ that ends the discussion.

\begin{remark}
The gain is smaller on real posteriors than on Gaussians ($2\times$, not $39\times$) because
neither law applies exactly to an anisotropic target. The direction is the same in every case
measured.
\end{remark}

\section{Obstruction 2: fixing the bandwidth buys time, and then a scale error}
\label{sec:fixed}

If the adaptive bandwidth is the problem, fix it. A large constant $h$ removes the feedback loop by
construction, and on an isotropic Gaussian a fixed $h \approx 100d$ --- two orders beyond the
$\sigma = \sqrt d$ that \cite{ba2022} tests and rejects --- has the correct fixed point $v = 1$.

On the real posteriors this works, and then does not. A fixed $h$ converts the collapse from an
attractor into a transient: the trajectory is a function of $h \times$ iterations, so ten times the
bandwidth buys ten times the delay and nothing else. Two starts a factor of $16$ apart in variance
meet at the same equilibrium and hold it to three digits for $450{,}000$ iterations --- so there
\emph{is} a fixed point, it is simply the wrong one.

\begin{remark}[a mechanism of ours that was wrong]
We first attributed the wrongness of that fixed point to anisotropy, on the grounds that a scalar
bandwidth cannot serve directions whose scales differ. That diagnosis survives as a description of
where the error goes, but the claim we built on it --- that there is no fixed point at fixed
bandwidth --- was an artifact of reading a single stiff-direction statistic. It is retracted.
\end{remark}

\subsection{The Laplace metric whitens these posteriors nearly perfectly}

If anisotropy is what is left, remove it. Whitening by $L$ with $LL^\top = H^{-1}$ at the mode
works better than we expected:

\begin{center}\small
\begin{tabular}{lrrl}
\toprule
system & $\operatorname{cond}(\Sigma_{\mathrm{ref}})$ & $\operatorname{cond}(L^{-1}\Sigma_{\mathrm{ref}}L^{-\top})$ & $5$--$95\%$ eigenvalues after\\
\midrule
fn     & $7.9\times10^{3}$ & \textbf{47}  & $0.86$ -- $1.22$\\
hiv    & $\mathbf{5.6\times10^{10}}$ & \textbf{1.6} & $0.82$ -- $1.21$\\
lorenz & $1.7\times10^{5}$ & 192 & $0.82$ -- $1.53$\\
\bottomrule
\end{tabular}
\end{center}

HIV's condition number falls by ten orders of magnitude. So preconditioning plus a fixed bandwidth
is the configuration to try, and it had never been run: earlier work tested preconditioning only at
the adaptive bandwidth, where it was the worst of four kernels.

\begin{remark}[a bug of ours]
That earlier verdict was worthless, and not because of the bandwidth. Our driver computed the
preconditioned score as \code{L.T @ jax.grad(lambda z: logp(x0 + L @ z))(y)}. Automatic
differentiation already differentiates through the \code{L @ z} inside, so it returns
$L^\top\nabla_x\log p$ on its own and the explicit \code{L.T} applied the metric a second time,
giving $L^\top L^\top \nabla_x \log p$. Verified against the analytic chain rule, the corrected line
agrees to $9.7\times10^{-16}$ and the old line is exactly the doubly-applied form to
$1.2\times10^{-15}$: a $98\%$ relative error. Every preconditioned measurement we had made was
invalid. None had been published.
\end{remark}

\subsection{With the bug fixed: an attractor, and a pure scale error}

Whitened, at $h = 10h^\star$ with $h^\star = 2d/\ln K$, and with a small \emph{fixed} step ---
Prodigy is unstable in whitened coordinates --- SVGD has a genuine attractor. The variance ratio in
the stiffest $10\%$ of directions, on FitzHugh--Nagumo at $K = 400$ over $100{,}000$ iterations:

\begin{center}\small
\begin{tabular}{lrrrr}
\toprule
start & $5$k & $25$k & $50$k & $100$k\\
\midrule
correct         & 0.959 & 0.787 & 0.738 & \textbf{0.748}\\
narrow $4\times$ & 0.077 & 0.137 & 0.244 & \textbf{0.494} $\uparrow$\\
wide $4\times$   & 13.18 & 4.703 & 1.587 & \textbf{0.788} $\downarrow$\\
\bottomrule
\end{tabular}
\end{center}

Approached from a factor of four below and a factor of four above. And the residual deficit is now
\emph{uniform}: $0.748$ against the $0.719$ the isotropic law predicts, a $1.04\times$ mismatch,
where the unpreconditioned dynamics leaves a $20\times$ mismatch between the mean and the stiff
directions. Preconditioning converts a shape error into a scale error.

A scale error is one number, and the Stein identity measures it without a reference. For a Gaussian
target and an ensemble of covariance $A$, $R = \tr(A\Sigma^{-1})/d$, so an ensemble uniformly a
factor $c$ too narrow reads $R = c$; rescaling about its own mean by $1/\sqrt R$ should correct it,
with no reference, no tuning and no extra gradient evaluations. It does:

\begin{table}[t]
\centering\small
\begin{tabular}{llrrrr}
\toprule
system & ensemble & Stein $R$ & energy $\times$ floor & stiff-var & max $|\theta$ err$|$\\
\midrule
fn     & \code{fit()} draws        & 1.274 & 1.65 & 1.058 & 0.077\\
       & SVGD equilibrium          & 0.702 & 2.42 & 0.748 & 0.242\\
       & \textbf{rescaled by $1/\sqrt R$} & 1.074 & \textbf{0.68} & \textbf{1.065} & 0.242\\
\midrule
lorenz & \code{fit()} draws        & 1.019 & 1.24 & 1.037 & 0.071\\
       & SVGD equilibrium          & 0.653 & 2.92 & 0.696 & 0.401\\
       & \textbf{rescaled by $1/\sqrt R$} & 1.074 & \textbf{0.96} & \textbf{1.066} & 0.401\\
\midrule
hiv    & \code{fit()} draws        & 1.000 & 1.03 & 1.074 & 0.096\\
       & SVGD equilibrium          & 0.538 & 6.78 & 0.573 & 0.067\\
       & \textbf{rescaled by $1/\sqrt R$} & 1.000 & \textbf{0.57} & \textbf{1.065} & \textbf{0.067}\\
\bottomrule
\end{tabular}
\caption{The reference-free recipe: start from \code{fit()} draws, whiten by $H^{-1}$, run at fixed
$h$ with a small fixed step to equilibrium, measure $R$, rescale by $1/\sqrt R$. A $2.4$ to
$6.8\times$ ensemble becomes a $0.57$ to $0.96\times$ one, from the ensemble alone. Note the last
column, which is the subject of Section~\ref{sec:theta}.}
\label{tab:recipe}
\end{table}

On the \emph{ensemble}, this is a success and we report it as one. On HIV it beats $400$ exact draws
on a Kolmogorov--Smirnov comparison. Obstruction 1 is removable.

\section{Obstruction 3: the fixed point is displaced in the parameters}
\label{sec:theta}

The last column of Table~\ref{tab:recipe} is the reason none of this is a recommendation. The
largest parameter error goes from $0.077$ to $0.242$ on FitzHugh--Nagumo and $0.071$ to $0.401$ on
Lorenz. The rescaling does not touch it, and cannot: it corrects a scale about the ensemble's own
mean, and this is an error \emph{in} that mean. Only HIV improves.

The energy distance improves while the parameters degrade because the energy distance is dominated
by the $322$ trajectory states. MAGI exists to estimate the $3$ to $7$ parameters.

\subsection{Where the displacement points}

It is not arbitrary. Decomposing the equilibrium's $\theta$ mean along the axis from the joint MAP
to the reference mean, with $t$ the signed fraction along that axis and \emph{gap} its length:

\begin{center}\small
\begin{tabular}{lrrrr}
\toprule
system & $t(\theta)$ & gap (sd) & $|t|\cdot$gap & measured $\theta$ error\\
\midrule
fn     & $+0.186$ & 1.029 & 0.191 & 0.242\\
lorenz & $-0.221$ & 1.803 & 0.399 & \textbf{0.401}\\
hiv    & $+0.454$ & 0.148 & 0.067 & \textbf{0.067}\\
\bottomrule
\end{tabular}
\end{center}

The displacement is aligned with that axis to $0.92$--$0.999$, against $0.58$ for a random
direction at $p = 3$, and two of the three products reproduce the measured error to the printed
digit. \textbf{HIV improves because its gap is $0.15$ posterior standard deviations, not because it
differs in kind.} The effect is specific to $\theta$: over the same run FitzHugh--Nagumo's state
block converges \emph{onto} the reference mean ($t$: $0.427 \to 0.024$) while $\theta$ leaves it.

\subsection{The mechanism is non-Gaussianity}

Run the identical dynamics --- same metric, bandwidth, step, start, particle count --- on an exact
Gaussian with the reference's own mean and covariance. The $\theta$ error lands at $0.0045$ sd on
FitzHugh--Nagumo and $0.0351$ on Lorenz, \emph{below} each reference's own Monte Carlo floor,
against $0.2422$ and $0.4013$ on the real posteriors.

So SVGD's fixed point is exact for a Gaussian target, and on a target that is not one it sits
displaced along the mode-to-mean axis by a fraction of that axis's length. That axis is not an
incidental feature of these problems. Correcting the mode to the mean is what the deterministic
pipeline is \emph{for} \cite{companion1}: the Laplace approximation at the MAP is mis-centred by a
full posterior standard deviation on a benchmark as benign as FitzHugh--Nagumo, and removing that
bias is the entire content of the method SVGD would be replacing. SVGD gives it back.

\begin{remark}[a second hypothesis of ours that was wrong]
The natural explanation is mode-seeking: the drift term is a kernel-weighted average of the score,
which should pull the ensemble mean toward high density. That predicts $t > 0$ everywhere. It is
wrong --- Lorenz's $t$ is $-0.221$, drifting \emph{away} from the MAP and past the reference mean
on the far side --- and the fraction is not universal either ($+0.186$, $-0.221$, $+0.454$). The
axis and the magnitude are understood; the sign is not.
\end{remark}

\subsection{What repairs it, and why that settles the question}

Importance-reweighting the equilibrium by the profiled log-density recovers parity with the
deterministic pipeline: $0.242 \to 0.071$, $0.401 \to 0.090$, $0.067 \to 0.029$, at effective
sample sizes of $78\%$, $66\%$ and $97\%$ and $\khat$ below $0.3$, in negligible time. On HIV that
is $3.3\times$ better than \code{fit()}.

This is the most useful negative result in the report. The repair works by reintroducing
$\log \hat p(\theta)$ --- the profiled density, which is the deterministic pipeline's entire
machinery. What does the estimating is the density; SVGD has supplied a proposal. And the
incumbent already makes its own proposal, from the profile mode and curvature, at a cost of four to
six seconds.

The drift is a finite-$K$ bias and does vanish, at $K^{-0.74}$ on FitzHugh--Nagumo and $K^{-0.36}$
on Lorenz. Reaching \code{fit()}'s accuracy needs roughly $1.5\times10^{4}$ and $1.2\times10^{5}$
particles respectively, and \code{fit()} is more accurate on $\theta$ at every $K$ we tested.

\begin{remark}[a third hypothesis of ours that was wrong]
We had called the step size a real lever and the $\mathrm{lr}\to 0$ limit the most consequential
untested gap. At matched flow time --- $\mathrm{lr}\times$iterations held constant --- three step
sizes spanning a factor of ten give $t = 0.186$ and a $\theta$ error of $0.2422$ identically to
four digits. The apparent dependence was less flow time. Withdrawn.
\end{remark}

\section{The one regime where dimension is not the obstacle}
\label{sec:profiled}

Obstruction 1 is about $d$. The natural response is to run SVGD where $d$ is small: on the profiled
$p$-dimensional marginal, with the states integrated out and autodiff taken through the inner
solve. At $p = 3$ this is the one configuration in which SVGD beat the incumbent --- $0.30\times$
and $0.36\times$ the $64$-draw floor on FitzHugh--Nagumo and Lorenz, better than \code{fit()} by
$2.5$ to $8\times$, at $6$ to $12\times$ the wall clock.

It does not survive $p = 5$. On HIV:

\begin{center}\small
\begin{tabular}{llrrrr}
\toprule
$K$ & method & $\theta$ energy & $\times$ floor & max $|\theta$ err$|$ & sd ratio\\
\midrule
16  & \code{fit()}   & 0.1193 & \textbf{0.50} & 0.119 & 0.961\\
16  & profiled SVGD  & 0.2529 & 1.05 & 0.103 & \textbf{0.578}\\
64  & \code{fit()}   & 0.0514 & \textbf{1.25} & 0.123 & 0.931\\
64  & profiled SVGD  & 0.0529 & 1.29 & 0.281 & \textbf{0.881}\\
256 & \code{fit()}   & 0.0145 & \textbf{0.76} & 0.103 & 0.986\\
\bottomrule
\end{tabular}
\end{center}

The sd ratios say why: the $\theta$ ensemble is itself underdispersed, which is \eqref{eq:law}
reappearing with $d$ replaced by $p$. It predicts $0.55$ at $(K=16, p=5)$ and $0.83$ at
$(K=64, p=5)$ against measured $0.33$ and $0.78$; at $p = 3$ it predicts $0.92$ at $K = 16$ and
above $1$ at $K = 64$, and FitzHugh--Nagumo and Lorenz duly succeed at every $K$. The law therefore
doubles as a design rule, and this is the third thing worth keeping. It is a rule about the
bandwidth as much as about the dimension, so it has two forms:

\begin{center}\small
\begin{tabular}{lrrrrrrr}
\toprule
$p$ & 2 & 3 & 4 & 5 & 6 & 7 & 8\\
\midrule
$K$ for $90\%$ of the correct variance, $h = \mathrm{Med}/\ln K$ & 96 & 192 & 512 & 2048 & 4096 & --- & ---\\
$K$ for $90\%$ of the correct variance, $h = \mathrm{Med}$ & \textbf{16} & \textbf{24} & \textbf{32} & \textbf{48} & \textbf{64} & \textbf{64} & \textbf{96}\\
\bottomrule
\end{tabular}
\end{center}

Measured on isotropic $\mathcal N(0, I_p)$, Adam at $0.05$, $4000$ iterations, started at exact
draws. Under the shipped $1/\ln K$ the requirement is $K \gtrsim 4e^{0.86p}$ --- $\ln
K_{\mathrm{crit}} = 0.862p + 1.449$, with the rate holding to $8$--$10\%$ over seven values of $p$
--- which is $320$ particles at $p=5$ and $2400$ at $p=7$. Dropping the log leaves the growth
exponential but slows the rate from $2.68\times$ to $1.33\times$ per dimension
($K_{\mathrm{crit}} \approx 10\,e^{0.287p}$), giving $48$ at $p=5$ and $64$ at $p=7$: a factor of
$43$ fewer particles at $p=5$, and the saving grows with $p$. So the $1/\ln K$ of
Section~\ref{sec:dim} is not merely a constant --- it changes how the cost scales.

\begin{remark}
The HIV measurements above were taken under the shipped $1/\ln K$, so they understate what
profiled SVGD can do: at $K = 64$, $p = 5$ the corrected bandwidth should reach about $0.93$ of the
correct variance rather than the $0.78$ measured. We have not re-run them, because the sd ratio is
not what decides this --- Section~\ref{sec:theta}'s displacement applies to the profiled marginal
too, and it is not a dispersion problem. The honest statement is that profiled SVGD's dispersion
at $p=5$ is better than these numbers show, and its parameter mean is not.
\end{remark}

Two further costs. HIV at $K = 256$ --- where the rule says it should work --- exhausted a $32$\,GB
V100 in the reverse pass, needing $256$ particles $\times$ three Cholesky factorisations of a
$603\times603$ block, and had to be rerun with the inner Newton steps rematerialised. And
Section~\ref{sec:theta}'s displacement applies here too, since the profiled marginal is no more
Gaussian than the joint posterior.

\section{Verdict}

Not worth pursuing, for the joint posterior or the profiled marginal.

The first obstruction alone would not settle it, because it is removable and we removed it. The
third does settle it, and it is structural rather than incidental: SVGD's fixed point is exact for
a Gaussian and displaced for anything else, along the one axis whose correction is the point of the
exercise. No bandwidth, metric or optimiser reaches it, because it is the fixed point and not an
approach to it; the only thing that repairs it reintroduces the density that makes SVGD
unnecessary. Against an incumbent that reaches the reference's own agreement floor on every
scorable system in four to six seconds and reports its own reliability without a reference, there
is nothing left to recommend.

Three things are worth keeping, and two of them are not about MAGI:
\begin{enumerate}[leftmargin=1.5em,itemsep=2pt]
\item \textbf{Drop the $1/\log K$} from the median-heuristic bandwidth. One line, $1.8$ to
  $2.4\times$ on real posteriors and up to $39\times$ on Gaussians, better in every case measured,
  and it replaces a rule nobody has analysed with one that has a theorem attached \cite{ba2022}.
\item \textbf{Report the Stein dispersion ratio.} $R \to 1$ under the target and $R < 1$ when
  underdispersed; for a Gaussian target it is exactly the mean variance-inflation factor, and in a
  smoke test it reads $0.9271$ where the measured ratio is $0.9270$. It is the only cheap
  diagnostic that sees SVGD's characteristic failure, which the optimiser's own convergence test
  cannot: the ensemble converges, to the wrong spread.
\item \textbf{A particle-count rule}, for anyone running SVGD on a low-dimensional marginal:
  $K \approx 10\,e^{0.287p}$ for $90\%$ of the correct variance under the plain median heuristic,
  or $4e^{0.86p}$ if the $1/\ln K$ is left in --- $48$ against $320$ particles at $p = 5$. Both
  are measured on isotropic Gaussians and optimistic for anything else.
\end{enumerate}

\section{What would change the verdict, and what we did not test}

We would revisit this if the displacement of Section~\ref{sec:theta} turned out to be removable by
something cheaper than the profiled density --- but the Gaussian control makes that unlikely, since
it shows the displacement is the fixed point's response to non-Gaussianity rather than an artifact
of finite $K$, finite time or the metric. The sign of $t$ is unexplained; explaining it would be
satisfying and would not change the magnitude.

Hes1 is untested throughout, and it is the one case where the incumbent is genuinely unreliable:
its gate declines, its effective sample size varies from $2.3\%$ to $21.1\%$ with the Sobol seed,
and its reference does not converge. It is also the only case where SVGD might have an argument, in
that it does not traverse the funnel geometry that defeats NUTS. We did not test it because with no
usable reference a success there could not be scored, and because Section~\ref{sec:theta}'s
mechanism is indifferent to the funnel. The better routes for Hes1 are elsewhere:
observe the unobserved component, reparameterise to the profiled coordinates, or Riemannian HMC
\cite{companion2}.

Finally, all measurements here are float64, at $K = 400$ unless stated, on two GPUs that agree to
four digits on a shared benchmark. Every reference-scored claim excludes Hes1.

\begin{thebibliography}{9}
\small
\bibitem{yang2021} S.~Yang, S.~W.~K.~Wong, S.~C.~Kou.
  Inference of dynamic systems from noisy and sparse data via manifold-constrained Gaussian
  processes. \emph{PNAS} 118(15), 2021.
\bibitem{liu2016} Q.~Liu, D.~Wang.
  Stein variational gradient descent: a general purpose Bayesian inference algorithm.
  \emph{NeurIPS}, 2016.
\bibitem{ba2022} J.~Ba, M.~A.~Erdogdu, M.~Ghassemi, S.~Sun, T.~Suzuki, D.~Wu, T.~Zhang.
  Understanding the variance collapse of SVGD in high dimensions. \emph{ICLR}, 2022.
\bibitem{huang2023} X.~Huang, H.~Dong, C.~Fang.
  Local KL convergence rate for Stein variational gradient descent with reweighted kernel, 2023.
\bibitem{companion1} A fast, self-certifying Gaussian posterior for MAGI ODE parameter inference.
  Companion technical report, 2026.
\bibitem{companion2} Profiling the states out: posterior inference for MAGI without a Gaussian
  assumption. Companion technical report, 2026.
\end{thebibliography}

\end{document}
